% =====================================================================
%  Seance 6 -- B. Cours 2 : perspectives (frames 11--17)
% =====================================================================
\section{Perspectives : ce qui vient après}

\begin{frame}{Perspectives 1 --- autoencodeurs pour métriques multivariées}
\small
\begin{columns}[t]
\begin{column}{0.52\textwidth}
\begin{block}{Le principe}
Un autoencodeur est un réseau de neurones entraîné à \textbf{reproduire ses
entrées} après un passage par un état compressé. Sur des données normales,
il reconstruit bien. Sur une anomalie, la reconstruction échoue : l'erreur
de reconstruction \textbf{est} le score d'anomalie.
\begin{itemize}
\item Avantage : capte les corrélations \textbf{entre} variables (CPU haut
+ latence normale $=$ suspect même si chaque variable est plausible).
\item Nos 14 variables de \file{nosys\_features.csv} deviennent un vecteur
d'entrée unique.
\end{itemize}
\end{block}
\end{column}
\begin{column}{0.44\textwidth}
\begin{wbox}[Pourquoi pas aujourd'hui]
Sur NOSYS, la forêt atteint déjà \textbf{AUC 0,80} avec 14 variables et
s'explique par importances. Un autoencodeur exige des milliers de fenêtres
normales, se règle difficilement, et son score \textbf{ne s'explique pas}.
Règle d'ingénieur : la complexité se gagne, elle ne s'emprunte pas.
\end{wbox}
\begin{ebox}[Quand il devient pertinent]
Quand le nombre de variables explose (centaines de métriques par service)
et que personne ne peut écrire les features à la main.
\end{ebox}
\end{column}
\end{columns}
\end{frame}

\begin{frame}{Perspectives 2 --- les LLM au NOC (culture, pas TP)}
\begin{center}\small
\begin{tabular}{@{}p{4.2cm}p{8.0cm}@{}}
\toprule
\textbf{Usage} & \textbf{Ce que ça change}\\
\midrule
Résumer un incident & à partir des logs de la nuit du 5/10 (brute force sur
\file{bdd-principal}, $\approx$ 125 échecs), produire en 10 secondes une
note : serveurs touchés, chronologie, actions posées\\
\addlinespace[1pt]
Requête en langage naturel & « donne-moi le p99 de latence de srv-01 la
semaine dernière » traduit en requête de métriques --- l'ingénieur ne
connaît plus la syntaxe par cœur\\
\addlinespace[1pt]
Rédaction de runbooks & transformer un historique de tickets en procédure
réutilisable\\
\bottomrule
\end{tabular}
\end{center}
\vspace{1mm}
\begin{wbox}[Les deux garde-fous à retenir]
\textbf{1.} Un LLM peut \textbf{se tromper avec assurance} : ses sorties sur
un incident se vérifient toujours contre les données brutes. \textbf{2.} On
n'envoie pas des logs contenant des noms d'utilisateurs ou d'adresses IP à
un service externe sans les pseudonymiser (on y revient en gouvernance).
\end{wbox}
\end{frame}

\begin{frame}{Perspectives 3 --- jumeaux numériques}
\small
\begin{columns}[t]
\begin{column}{0.52\textwidth}
\begin{block}{Le principe}
Un \textbf{jumeau numérique} est une simulation du système d'information
assez fidèle pour qu'on y teste des scénarios sans toucher la production :
\begin{itemize}
\item rejouer l'incident du 22/09 et vérifier que le modèle l'aurait
détecté \textbf{plus tôt} ;
\item tester un nouveau seuil ou un nouveau modèle sur \textbf{des mois
d'incidents simulés} en quelques minutes ;
\item entraîner l'équipe du NOC sur des pannes qui n'ont jamais eu lieu.
\end{itemize}
\end{block}
\end{column}
\begin{column}{0.44\textwidth}
\begin{ebox}[Le lien avec nos données]
Nos 30 jours de métriques NOSYS (graine 5534) sont déjà une
\textbf{réplique simplifiée} : un générateur de charge + un injecteur de
défauts. Un jumeau, c'est la même idée avec la topologie réseau en plus.
\end{ebox}
\begin{wbox}[La limite honnête]
Un jumeau reproduit les pannes \textbf{qu'on sait décrire}. La panne
véritablement nouvelle (le cygne noir) n'existe que dans la production ---
d'où le monitoring de la séance précédente, qui reste irremplaçable.
\end{wbox}
\end{column}
\end{columns}
\end{frame}

\begin{frame}{Perspectives 4 --- edge ML : détecter sur l'équipement}
\small
\begin{columns}[t]
\begin{column}{0.52\textwidth}
\begin{block}{Le déplacement du calcul}
Aujourd'hui : les agents envoient les métriques à un serveur central qui
score (notre pipeline). En \textbf{edge ML}, le modèle tourne
\textbf{sur l'équipement} lui-même --- routeur, serveur, passerelle IoT.
\begin{itemize}
\item Détection locale en millisecondes, même si le réseau est coupé
(exactement le scénario d'une vraie panne).
\item Moins de données qui circulent : les métriques brutes restent locales,
seules les alertes partent.
\end{itemize}
\end{block}
\end{column}
\begin{column}{0.44\textwidth}
\begin{wbox}[Le contrainte : la latence, encore]
Sur un équipement, pas de 200 arbres à 60 ms : il faut des modèles
\textbf{compilés et quantifiés} (par ex. arbre unique, modèle linéaire,
micro-réseaux). Le TP B2 prend son sens : mesurer la latence AVANT de
décider où le modèle peut vivre.
\end{wbox}
\begin{kbox}[Fil conducteur]
Toutes ces perspectives (autoencodeurs, edge, jumeaux) changent le
\textbf{modèle ou son lieu d'exécution} --- jamais les fondations :
données propres, split temporel, validation, monitoring.
\end{kbox}
\end{column}
\end{columns}
\end{frame}

\begin{frame}{Quiz 2 --- perspectives (5 min)}
\begin{qbox}[3 questions]
\begin{enumerate}
\item Le score d'anomalie d'un autoencodeur est :\\
(a) la sortie du réseau \quad (b) l'erreur de reconstruction \quad (c) la
distance entre deux points
\item Pourquoi un autoencodeur est-il difficile à justifier devant un
exploitant ?\\
(a) il est trop rapide \quad (b) son score ne s'explique pas par des
variables nommées \quad (c) il exige des étiquettes
\item Le principal intérêt du edge ML pour la détection de panne :\\
(a) plus de précision \quad (b) détecter même si le lien réseau est coupé
\quad (c) un modèle plus gros
\end{enumerate}
\end{qbox}
\vspace{1mm}
\begin{solbox}[Réponses]
\textbf{1 : b} \quad \textbf{2 : b} --- l'interprétabilité reste un critère
de production \quad \textbf{3 : b}
\end{solbox}
\end{frame}
